\chapter{Anexos}

\section{Estructura del repositorio}

Carpetas relevantes:
\begin{itemize}[leftmargin=1.5cm]
  \item \texttt{src/algoritmos}
  \item \texttt{src/estructuras}
  \item \texttt{src/universidad-proyecto} (datos, lógica, API, frontend)
  \item \texttt{src/docs/sections} (archivos \texttt{.tex} de esta documentación)
\end{itemize}

\section{Datasets}

Archivos incluidos: \texttt{src/universidad-proyecto/data/data.json} y \texttt{data2.json}. Uno de ellos se usa como instancia principal para la validación y el otro queda como entrada alternativa.

\section{Imágenes}

Nombres esperados en \texttt{figuras/}:
\begin{itemize}[leftmargin=1.5cm]
  \item \texttt{frontend-horario-global.png}
  \item \texttt{frontend-formulario-estudiante.png}
  \item \texttt{frontend-horario-estudiante.png}
  \item \texttt{flujo-algoritmos.png}
\end{itemize}

\section{Comandos útiles}

Ejemplo de arranque del servidor:
\begin{verbatim}
npx ts-node src/universidad-proyecto/index.ts
\end{verbatim}
